% Source: https://artofproblemsolving.com/community/c3h1443851_vietas_formulas_question
% Stable identity: aops-topic-1443851-q01
\textbf{(a)} Let $t=\sqrt[3]{7}$, so $t^3=7$ and $x=t^2+t$. Then
\[
x^3=(t^2+t)^3=t^6+3t^5+3t^4+t^3=49+21t^2+21t+7=56+21x .
\]
So $x$ is a root of $\boxed{p(y)=y^3-21y-56}$.

\textbf{(b)} Let $\omega$ be a non-real cube root of $1$, so $\omega^2+\omega+1=0$ and $\omega^3=1$. The computation in (a) used only $t^3=7$, so it works for $t=\omega\sqrt[3]{7}$ and $t=\omega^2\sqrt[3]{7}$ as well. Hence the three roots of $p$ are
\[
t^2+t\quad\text{for}\quad t=\sqrt[3]7,\ \omega\sqrt[3]7,\ \omega^2\sqrt[3]7,
\]
that is, $\sqrt[3]{49}+\sqrt[3]{7}$, $\ \omega^2\sqrt[3]{49}+\omega\sqrt[3]{7}$ and $\ \omega\sqrt[3]{49}+\omega^2\sqrt[3]{7}$. They are distinct, because $t\mapsto t^2+t$ takes different values at these three numbers (the difference of two values is $(t-u)(t+u+1)$, and $t+u+1\neq0$ since $|t+u|=\sqrt[3]7>1$). The last two are complex conjugates and are not real, with real part $-\dfrac{\sqrt[3]{49}+\sqrt[3]{7}}{2}$, so $p$ has only one real root.

So, unlike square roots, the ``conjugate'' $\sqrt[3]{49}-\sqrt[3]{7}$ is \textbf{not} a root: it is real and different from $\sqrt[3]{49}+\sqrt[3]7$. Numerically $p\!\left(\sqrt[3]{49}-\sqrt[3]{7}\right)\approx 1.74^3-21(1.74)-56\approx-87\neq0$. For cube roots the role of the sign change $\sqrt b\mapsto-\sqrt b$ is played by multiplying the cube root by $\omega$ or $\omega^2$.
